%---------------------------------------------------------------
% CHAPTER 1: History of data centers
\chapter{Historie DC a dnešní nároky}
\label{chap:historie}

Digitální svět je často vnímán jako nehmotná entita existující v cloudu, ale jeho fyzické základy v podobě datových center procházejí nejvýraznější transformací od dob prvních sálových počítačů. Paradigma datového centra se totiž přesunulo od skladu informací k superpočítači provádějícímu náročné výpočty.~\cite[s.~11]{ranganathan_twenty_2024}

Právě tento technologický skok, doprovázený přechodem ke kapalinovému chlazení~\cite{bizo_silicon_2022}, definuje nové nároky na infrastrukturu i její okolí. Následující text analyzuje historický vývoj těchto staveb, technologické příčiny jejich rostoucí náročnosti a faktory, které v současnosti určují umístění datových center.

\section{Historie a architektura}
\label{sec:historie}

%%%%%%%%%%%%%%%%%%%%%%%%%%%%%%%%%%%%%%%%%%%%%%%%%%%%%%%%%%%%%%%%%%%%%%%%%%%%%
% SEKCE 1: HISTORIE A ARCHITEKTURA
% - Definice datového centra napříč historií
% - Uvedení klasifikačního rámce Tier od Uptime Institute
% - Definice WSC a důvody přechodu od tradičních DC
% - Popis moderního AI datového centra + AI factory]
%%%%%%%%%%%%%%%%%%%%%%%%%%%%%%%%%%%%%%%%%%%%%%%%%%%%%%%%%%%%%%%%%%%%%%%%%%%%%

Historicky byla datová centra chápána jako zázemí pro ukládání dat a distribuované výpočty, nicméně s rozvojem cloudu a následným nástupem specializovaných AI továren se tato role zásadně proměnila. Tato sekce popisuje, jak se z původních ubytoven pro servery staly integrované výpočetní celky, které dnes fungují jako jeden monolitický superpočítač. \cite{ranganathan_twenty_2024, barroso_data_2026, barroso_brief_2021}

\subsection{Tradiční pojetí DC a standardy spolehlivosti}
\label{subsec:tradice}

Rihards Balodis a Inara Opmane~\cite[s.~180--181]{tatnall_reflections_2012} v roce 2012 definovali datové centrum jako specializované fyzické prostředí určené pro bezpečné umístění počítačových systémů a jejich komponent. Jeho nezbytné zázemí tvoří systémy pro distribuci a zálohování elektřiny, redundantní komunikační rozvody, klimatizace, systémy požární ochrany a prvky fyzického zabezpečení objektu.

Modernější definice datových center zahrnují i samotná výpočetní zařízení. Například podle odborného článku z roku 2020 v časopise Facilities~\cite{fadaeefath_abadi_data_2020} je DC definováno jako významná infrastruktura obsahující velké množství serverů a počítačů, která poskytuje internetové služby společnostem po celém světě.

Historicky tak byla datová centra vnímána především jako \textit{budovy}, jež mají poskytovat vhodné podmínky pro provoz počítačových systémů. Novější definice vnímá datové centrum nejen jako budovu, ale jako soubor vybavení pro zpracování, ukládání a přeposílání informací. Oba texty se však shodují v tom, že datová centra vyžadují speciální komunikační nástroje a subsystémy pro napájení, regulaci okolní teploty a vlhkosti vzduchu~\cite{tatnall_reflections_2012, fadaeefath_abadi_data_2020}.

Výše uvedené klíčové subsystémy ovlivňují výslednou dostupnost služeb pro koncové uživatele. Za účelem sjednocení těchto požadavků navrhla organizace Uptime Institute čtyři třídy DC v rámci celosvětově uznávané klasifikace Tier~\cite{turner_iv_tier_2008}. Tato certifikace definuje úrovně, které se od sebe liší zejména mírou redundance komponent a konfigurací distribučních cest pro provoz IT vybavení.

Souhrn parametrů pro jednotlivé úrovně Tier uvádí tabulka~\ref{tab:uptime_tiers}, detailnější popis je obsažen v příloze~\ref{appendix:tiery}. Tato klasifikace tvoří vrchol éry, v níž byla spolehlivost definována fyzickou redundancí komponent. S nástupem virtualizace se však pozornost přesunula k softwarově definované infrastruktuře.

\begin{table}[t] \centering \setlength{\tabcolsep}{4pt}

    \caption[Souhrn požadavků pro klasifikaci Tier]{Souhrn požadavků pro klasifikaci Tier~\cite[s.~13]{turner_iv_tier_2008}. N označuje kritickou kapacitu nezbytnou pro podporu aktuálního IT zatížení}
    \label{tab:uptime_tiers}
    
    \begin{tabular}{lcccc}
        \toprule
        \textbf{Parametr} & \textbf{Tier I} & \textbf{Tier II} & \textbf{Tier III} & \textbf{Tier IV} \\
        \midrule
        Dostupnost & 99,671~\% & 99,741~\% & 99,982~\% & 99,995~\% \\
        Redundance komponent & N & N+1 & N+1 & >2N \\
        Distribuční cesty & 1 & 1 & 1 akt., 1 zálož. & 2 akt. \\
        Souběžná udržovatelnost & Ne & Ne & Ano & Ano \\
        Odolnost proti poruchám & Ne & Ne & Ne & Ano \\
        \bottomrule
    \end{tabular}

\end{table}

\subsection{Koncept integrovaného systému WSC}
\label{subsec:wsc}

Na přelomu tisíciletí se začalo formovat nové paradigma distribuovaných výpočtů. Namísto delegování zátěže na koncová zařízení (osobní počítače) se požadavky na výkon a kapacitu přesunuly do centralizované infrastruktury~\cite{barroso_data_2026}. Tento koncept poskytování škálovatelného výpočetního výkonu, softwaru a úložiště prostřednictvím internetu se nazývá Cloud Computing~\cite{lewis_basics_2010}.

Právě Cloud Computing umožnil vývoj softwarových systémů s nároky, které řádově přesahovaly možnosti izolovaných uživatelských zařízení -- úzkým místem se tak stala infrastruktura samotného provozovatele softwaru. V reakci na tyto potřeby vznikly mohutnější a výkonnější výpočetní celky s vyšší kapacitou, které L. A. Barroso definuje jako Warehouse-Scale Computers (WSC)~\cite{barroso_data_2026, barroso_brief_2021}.

Barroso rozdíl mezi WSC a tradičním DC ilustruje v pěti dimenzích hardwarové a softwarové architektury~\cite[s.~2--4]{barroso_data_2026}:

\begin{description}
    \item[Homogenita] Tradiční datová centra hostí různorodý hardware mnoha firem. WSC zpravidla využívá jednotný hardware pod kontrolou jediné organizace, aby byla zajištěna rychlá nahraditelnost a nákladová efektivita.
    
    \item[Software] Zatímco běžný poskytovatel dodává pouze budovu, konektivitu a fyzickou správu hardwaru~\cite[s.~10]{marx_gomez_foundations_2017}, WSC disponuje komplexní softwarovou vrstvou, která propojuje tisíce serverů a umožňuje jim komunikovat.
    
    \item[Škálovatelnost] WSC je navržen pro běh masivních internetových služeb využívajících desetitisíců serverů současně. To je umožněno díky efektivní komunikaci mezi servery, kterou běžná DC neposkytují.
    
    \item[Odolnost] Tradiční centra se snaží chybám hardwaru předcházet, což je nákladné. WSC naopak s výpadky komponent počítá a řeší je pomocí architektury navržené pro automatickou (graceful) toleranci jejich poruch.
    
    \item[Měřítko] WSC bývají o jeden až dva řády větší než běžná podniková datová centra, přičemž veškeré prostředky (dodávky energie, chlazení) jsou sdíleny v rámci společné infrastruktury pro správu zdrojů.
\end{description}

Technologický vývoj se mezi lety 2004 a 2008 soustředil na optimalizaci provozu. Důležitou inovací bylo zavedení metriky PUE, která je popsána v podsekci~\ref{subsec:pue}. Později se začaly objevovat první úlohy zaměřené na hluboké učení (např. projekt Google Brain)~\cite{ranganathan_twenty_2024}, což předznamenalo přerod univerzálních WSC v úzce specializovaná AI centra.

\subsection{Architektura moderních AI továren}
\label{subsec:ai_factory}

Nejmodernější druh datových center se nazývá AI Factory (továrna na AI). Podle slovníku společnosti NVIDIA, předního dodavatele hardwaru pro AI datová centra, jde o specializovanou výpočetní infrastrukturu, která vytváří hodnotu spravováním dat po celou dobu životního cyklu AI.~\cite{nvidia_what_2026}

Popisovaný životní cyklus je iterativní proces zahrnující sbírání dat, trénování a ladění modelu až po jeho hodnocení. Poté se model přesouvá do fáze \textit{inference}, tedy přeměny vstupu na výstupy (dotazování modelu). Koncept AI továrny představuje moderní přístup, který se snaží tyto fáze propojit v jednom datovém centru. V praxi se však zatím provozují zpravidla odděleně:

\begin{itemize}
    \item \textbf{Trénovací centra} jsou navržena pro vysoký a konzistentní výkon, pročež vyžadují pokročilé metody chlazení~\cite{cruzes_data_2025}. Komunikace probíhá v rámci jednoho DC, takže mohou být v odlehlých lokalitách. Jde o novou technologii.

    \item \textbf{Inferenční centra} kladou důraz na efektivitu a nízkou latenci, kvůli čemuž by měla být stavěna v blízkosti optických sítí. Jejich zátěž je vysoce stochastická~\cite{cruzes_data_2025, ginzburg-ganz_technical_2025}. Podobají se WSC a jsou již dobře zavedená.
\end{itemize}

Rozdíl mezi trénovacími centry a WSC spočívá v přechodu od izolovaných úloh k realizaci masivně paralelizovaných výpočtů. Barrosova definice WSC počítá s nezávislostí jednotlivých serverů (jakýkoli server lze díky virtualizaci nahradit). Naproti tomu TC funguje jako jeden logický superpočítač, takže porucha jediného procesoru může zastavit trénování celého modelu~\cite[s.~270]{barroso_data_2026}. Tato provázanost definuje trénovací centrum jako úzce specializovaný stroj, kde je integrace a rychlost propojení důležitější než kdykoli předtím.

Dalším rozdílem je požadovaná dostupnost: Zatímco tradiční datová a inferenční centra bývají zpravidla Tier III nebo IV, trénovací centra mohou využívat architekturu s nižší úrovní fyzické redundance, často odpovídající úrovni Tier I. Případný výpadek pro ně díky postupnému ukládání stavu nepředstavuje kritické ohrožení služby, ale pouze dočasnou ztrátu výpočetního času~\cite[s.~83]{barroso_data_2026}. Tato optimalizace pro nákladovou efektivitu a rychlost výstavby na úkor drahé dostupnosti umožňuje zkrátit čas potřebný pro vývoj nových modelů, což je pro rychle se vyvíjející odvětví z byznysového hlediska důležité~\cite[s.~24]{ginzburg-ganz_technical_2025}.

\section{Konstrukce a hardware}
\label{sec:konstrukce}

%%%%%%%%%%%%%%%%%%%%%%%%%%%%%%%%%%%%%%%%%%%%%%%%%%%%%%%%%%%%%%%%
% SEKCE 2: KONSTRUKCE A HARDWARE
% - Vnitřnosti DC - racky, servery, procesory, switche, ...
% - Přechod ke specializovaným čipům, CPU -> GPU -> TPU
% - Uspořádání v datovém centru, co obvykle zahrnují z prostorového hlediska
%%%%%%%%%%%%%%%%%%%%%%%%%%%%%%%%%%%%%%%%%%%%%%%%%%%%%%%%%%%%%%%%

Následující část se zaměřuje na jednotlivé komponenty tvořící a podporující výpočetní výkon datových center. Stěžejním tématem je standardizace serverových skříní, fyzická konstrukce moderních akcelerátorů a členění vnitřních prostorů datového centra, které musí splňovat extrémní nároky na statické zatížení budov. Zdrojům energie a chlazení DC se věnují sekce \ref{sec:energie} a \ref{sec:chlazeni}.

\subsection{Fyzické parametry serverových skříní}
\label{subsec:racky}

Základním stavebním kamenem každého datového centra je \textit{rack}. Jde o standardizovaný kovový rám, do kterého se horizontálně zasouvají a upevňují modulární komponenty (servery, úložiště, síťové přepínače, směrovače, \dots{}). Vnější konstrukce standardního racku je přibližně 1 980~mm vysoká, 580--640~mm široká a 660--760~mm hluboká, ale existují i jiné velikosti.~\cite[s.~2940]{kant_data_2009}

Podle standardu EIA-310-D má mít běžný rack montážní šířku 19 palců (482,6~mm), což odpovídá standardní šířce předního panelu IT vybavení. Stejný standard definuje jednotku U (1~U = 44,45~mm), která představuje montážní výšku typického IT zařízení. Pokud je tedy rack popsán jako 42U, dokáže pojmout 42 zařízení s výškou 1~U. Specifické komponenty, zejména výkonné AI akcelerátory, však mohou mít kvůli chlazení výšku 4~U, 8~U nebo i více.~\cite{hu_how_2020}

Z těchto důvodů standardní racky pro potřeby výkonných AI serverů nepostačují. V. Avelar~\cite[s.~13--14]{avelar_ai_2023} doporučuje, aby byly racky pro trénování umělé inteligence širší (např. 750~mm místo 600~mm pro uložení rozvodů kapalinového chlazení), hlubší (přes 1200~mm) a vyšší (48~U a více). Plně osazený AI rack může vážit přes 900~kg, což je potřeba zvážit při jeho přepravě i návrhu DC.

\subsection{Čipy a specializované akcelerátory}
\label{subsec:hardware}

Historicky byla datová centra poháněna univerzálními procesory (CPU), které jsou vynikající pro sériové zpracování složitého, větveného a nepravidelného kódu. Trénování umělé inteligence je však výpočetně velmi specifické -- spoléhá na několik málo operací, které jsou paralelně aplikovány na obrovské množství dat~\cite[s.~133--134]{barroso_data_2026}. CPU tak byly nahrazeny jinými druhy čipů, které jsou pro takto specifické úkoly efektivnější~\cite[s.~1391]{phani_suresh_paladugu_evolution_2025}.

Mezi používané komponenty patří především grafické procesory (GPU), které obsahují tisíce jednodušších jader (architektura SIMT) užitečných při paralelizaci. Dalším evolučním krokem jsou čipy navržené čistě pro potřeby AI, což je činí efektivnějšími. Za zmínku stojí TPU (Tensor Processing Units) vyvinuté společností Google, které využívají systolická pole -- propojenou síť výpočetních jednotek optimalizovaných pro maticové násobení.~\cite[s.~134--139]{barroso_data_2026}

Některé analogie uvádějí, že CPU je jako sportovní auto, protože dokáže velice rychle přepravit pár lidí. Naproti tomu GPU jsou jako autobusy, které sice jedou pomaleji a nevejdou se na všechny cesty, ale dokážou přepravit stovky lidí najednou. Nejsou kvalitativně horší ani lepší, jen slouží k jinému účelu.

\subsection{Vnitřní uspořádání a technické zóny}
\label{subsec:dispozice}

Vnitřní dispozice datového centra se člení na dvě hlavní oblasti: White Space a Gray Space. První představuje plochu vyhrazenou pro IT vybavení (racky, síťové prvky a servery). Gray Space naopak zahrnuje plochu pro veškerou podpůrnou infrastrukturu (transformátory, UPS a systémy chlazení)~\cite[s.~300]{geng_data_2014}. U tradičních DC tvořilo technické zázemí menší část celkové dispozice, přičemž poměr mezi Gray a White Space byl u větších datových center až 1:3~\cite{butler_unlocking_2025}.

Extrémní výkonová hustota AI racků však tento poměr mění: Jednak zmenšuje nároky na podlahovou plochu v rámci počítačového sálu (výkon je mnohem více koncentrovaný), jednak vyžaduje větší napájecí a chladicí systémy, které fyzicky zabírají stále více prostoru (zejména kvůli mohutnému potrubí). Jejich vzájemný poměr se tak může posunout až nad hranici 1:1~\cite{butler_unlocking_2025}.

\section{Energie}
\label{sec:energie}

%%%%%%%%%%%%%%%%%%%%%%%%%%%%%%%%%%%%%%%%%%%%%%%%%%%%%%%%%%%%%%%%
% SEKCE 3: ENERGIE
% - Co je to příkon a jaký je typický příkon jednoho datového centra
% - Jak se měří efektivita DC -- PUE
% - Odkud datová centra berou energii (důležité pro další kapitoly)
%%%%%%%%%%%%%%%%%%%%%%%%%%%%%%%%%%%%%%%%%%%%%%%%%%%%%%%%%%%%%%%%

Energetická infrastruktura se stala hlavním limitujícím faktorem moderních datových center. Tato kapitola analyzuje, jakým způsobem revoluce v oblasti umělé inteligence narušila dekádu trvající stabilitu energetické spotřeby a jaké nároky klade nejen na samotná DC, ale také na sdílenou energetickou síť.

\subsection{Příkon a spotřeba datových center}
\label{subsec:prikon}

Při analýze energetické náročnosti je nezbytné rozlišovat mezi příkonem (jednotka watt) a spotřebou (jednotka watthodina). V tomto kontextu příkon představuje okamžitý odběr elektřiny, který v čase kolísá podle vytížení DC. Příkon dále rozdělujeme na dvě kategorie: celkový příkon definuje horní hranici odběru celého datového centra při 100\% vytížení, zatímco instalovaný příkon (IT příkon) definuje horní hranici odběru samotného IT vybavení. Spotřeba pak vyjadřuje množství elektřiny, která byla odebrána za určitý čas.

Zatímco menší podniková datová centra měla historicky příkon menší než 1~MW~\cite{barroso_data_2026, sheng_power_2026}, větší datová centra určená pro Cloud Computing mají příkon 10 až 50~MW~\cite[s.~38]{international_energy_agency_iea_energy_2025}. Moderní hyperscale AI trénovací klastry však fungují v úplně jiných mezích -- jedno takové zařízení vyžaduje příkon 100~MW až 1 GW~\cite{sheng_power_2026, chen_electricity_2025}. Přehled instalovaného příkonu datových center je uveden v tabulce~\ref{tab:energeticke_rozsahy}.

\begin{table}[t] \centering

    \caption[Srovnání instalovaného příkonu]{Srovnání instalovaného příkonu dle typu datového centra. Vlastní zpracování podle dat od Barrosa~\cite{barroso_data_2026}, IEA~\cite{international_energy_agency_iea_energy_2025} a Chen~\cite{chen_electricity_2025}}
    \label{tab:energeticke_rozsahy}

    \begin{tabular}{lcl}
        \toprule
        \textbf{Typ DC} & \textbf{Instal. příkon} & \textbf{Zaměření} \\
        \midrule
        Enterprise DC & < 1~MW & Lokální firemní IT \\
        Cloud Computing & 10--50~MW & SaaS, IaaS \\
        Hyperscale AI & 100--1~000+~MW & Trénování LLM \\
        \bottomrule
    \end{tabular}

\end{table}

Instalovaný příkon není směrodatný pro spotřebu -- je nutné zvážit jeho skutečné využití. Cloudová DC jsou typicky využívána z asi 33,5~\%~\cite[s.~224]{barroso_data_2026}, u AI trénovacích center se využití blíží 100~\%. Jinými slovy, po dobu trénování modelu se průměrný příkon blíží instalovanému příkonu~\cite[s.~3]{avelar_ai_2023}. Příkladem: Pokud ve 200MW datovém centru běží AI servery 80~\% času~\cite[s.~27]{LBNL-2001637} (zbylých 20~\% jsou prostoje), jeho roční spotřeba činí asi 1,401~TWh. To odpovídá spotřebě Karlovarského kraje za rok 2024 ($\approx$1,5~TWh)~\cite{energeticky_regulacni_urad_spotreba_2024}.

\subsection{Power Usage Effectiveness}
\label{subsec:pue}

Pro zvýšení efektivity je nutné co největší část spotřebované energie využít na samotný výpočetní výkon, což ale není možné bez podpůrné infrastruktury. Z tohoto důvodu je nutné podrobně analyzovat, kam a v jakém množství energie skutečně proudí. Spotřebu energie v DC proto dále rozdělujeme na~\cite{sheng_power_2026, chen_electricity_2025}:

\begin{itemize}
    \item \textbf{IT vybavení} (servery, disková pole, sítě): zhruba 40--50~\% spotřeby.

    \item \textbf{Chladicí infrastrukturu}: tradičně spotřebovává 30--40~\% dodané energie.

    \item \textbf{Podpůrná zařízení} (napájení, osvětlení, zabezpečení): zbylých 10--30~\%.
\end{itemize}
 
K měření celkové efektivity se používá metrika PUE (Power Usage Effectiveness), která udává poměr mezi celkovou energií dodanou do budovy a energií, která skutečně dorazí k IT vybavení~\cite{chen_electricity_2025}:

$$\text{PUE} = \frac{\text{Total Facility Electricity Consumption}}{\text{IT Equipment Electricity Consumption}} \geq 1.$$

Zatímco v roce 2010 byla průměrná hodnota PUE datových center po celém světě kolem 2,5 (tzn. na každý 1 watt pro server připadlo 1,5 wattu na chlazení a ztráty), průměr se dnes podařilo snížit na 1,55~\cite{de_roucy-rochegonde_ai_2025}. Moderní AI továrny však míří k vyšší efektivitě -- Google a Meta uvádějí, že jejich vybraná DC mají PUE pod 1,1~\cite{chen_electricity_2025, sheng_power_2026}, pravděpodobně díky efektivnímu chlazení vodou.

\subsection{Odkud datová centra berou energii}
\label{subsec:odkud_energie}

Většina datových center je napájena z veřejné elektrické sítě. Hyperscale centra se často připojují k vysokonapěťové soustavě (v USA 115--230~kV), zatímco menší DC využívají střední napětí~\cite[s.~4]{chen_electricity_2025}. Původ této energie tak závisí na energetickém mixu daného regionu, přičemž celosvětově v napájení DC převažují fosilní paliva. Podíl OZE má v budoucnu růst.~\cite[s.~86]{international_energy_agency_iea_energy_2025}

Kromě toho je každé centrum vybaveno vlastními záložními (naftovými, plynovými) generátory nebo velkokapacitními palivovými články. Ty mohou sloužit nejen jako záloha, ale i jako primární napájení. Obrovská spotřeba podněcuje trend umisťování AI DC přímo ke zdrojům energie: velké firmy je začínají stavět v bezprostřední blízkosti jaderných elektráren.~\cite{chen_electricity_2025, cruzes_data_2025, international_energy_agency_iea_energy_2025} 

\section{Voda a chlazení}
\label{sec:chlazeni}

%%%%%%%%%%%%%%%%%%%%%%%%%%%%%%%%%%%%%%%%%%%%%%%%%%%%%%%%%%%%%%%%
% SEKCE 4: CHLAZENÍ
% - TDP jakožto ukazatel zvyšujících se nároků na chlazení
% - Nutnost přechodu od vzduchového ke kapalinovému chlazení
% - Dopady na spotřebu a odběr vody DC
% - TODO: Snažit se zakomponovat dodatky do normálních kapitol
%%%%%%%%%%%%%%%%%%%%%%%%%%%%%%%%%%%%%%%%%%%%%%%%%%%%%%%%%%%%%%%%

Výkonné čipy není možné chladit vzduchem kvůli jeho fyzikálním vlastnostem, pročež byli designéři datových center nuceni přejít na kapalinové chlazení. Voda totiž disponuje mnohem vyšší tepelnou kapacitou a vodivostí než vzduch, takže dokáže odvést více tepla za kratší čas. To však představuje nový problém -- DC se stala závislá na dalším sdíleném zdroji, který musí odněkud čerpat.

\subsection{Thermal Design Power}
\label{sec:tdp}

Metrika TDP (Thermal Design Power) určuje maximální tepelný výkon jednoho čipu při běžném provozu~\cite{mahajan_cooling_2006}, což je důležité pro návrh chladicího systému. Z prvního zákona termodynamiky zároveň plyne, že TDP je přibližně roven doporučenému příkonu čipu~\cite{avelar_ai_2023}. Tato metrika však neurčuje \textit{maximální} generované teplo -- to může být při intenzivních výpočtech krátkodobě vyšší~\cite{gigabyte_tdp_2026}.

Ze získané literatury plyne, že za posledních 10 let došlo u IT komponent k obrovskému nárůstu TDP, který způsobily především moderní AI čipy. Nejvýkonnější GPU dnes dosahují TDP až 1400~W, historický přehled je v tabulce~\ref{tab:gpu_tdp}. 

\begin{table}[b] \centering

    \caption[Přehled vývoje TDP]{Přehled vývoje TDP grafických procesorů NVIDIA. Tabulka převzata od Schneider Electric~\cite{avelar_ai_2023} a doplněna o nové údaje z dokumentace NVIDIA~\cite{nvidia_nvidia_2025, nvidia_nvidia_2025-1}. Srovnání: CPU servery měly v roce 2010 TDP kolem 120~W, dnes mají asi 400~W~\cite{bizo_silicon_2022}}
    \label{tab:gpu_tdp}

    \begin{tabular}{lccc}
        \toprule
        \textbf{GPU} & \textbf{Paměť} & \textbf{Rok vydání} & \textbf{TDP} \\
        \midrule
         V100 SXM2 & 32~GB & 2018 & 300~W \\
         A100 SXM & 80~GB & 2020 & 400~W \\
         H100 SXM & 80~GB & 2022 & 700~W \\
         B200 (Blackwell) & 186~GB & 2025 & 1000~W \\
         B300 (Blackwell Ultra) & 279~GB & 2025 & 1400~W \\
        \bottomrule
    \end{tabular}

\end{table}

Mnohem větší problém než růst TDP však představuje fyzická blízkost čipů, která je vynucena požadavky na rychlou synchronizaci při trénování modelů. Na větší vzdálenosti se totiž musí použít velmi drahá~\cite[s.~4]{avelar_ai_2023} a energeticky náročná~\cite[s.~36]{cruzes_data_2025} optická vlákna, která synchronizaci i tak zpomalují. Návrháři komponent se proto snaží umístit co nejvíce čipů do jednoho serveru, což vede k prostorové koncentraci tepla. Nejnovější AI racky dosahují hustoty výkonu přes 100~kW (běžně 30 až 80~kW), která z fyzikálních důvodů přesahuje možnosti vzduchového chlazení a vynucuje si přechod na chlazení kapalinové.~\cite[s.~23]{cruzes_data_2025}

\subsection{Od vzduchového chlazení ke kapalinovému}
\label{subsec:vzduch}

Tradiční chlazení vzduchem historicky dominovalo díky své jednoduchosti a nízké ceně, ale postačuje pro racky se spotřebou do 20~kW~\cite[s.~10]{avelar_ai_2023}. U AI racků spotřeba stoupá až na 100~kW, kvůli čemuž se kapalinové chlazení stalo absolutní nutností -- kapaliny teplo odvádějí efektivněji než vzduch~\cite[s.~4]{chen_electricity_2025}. Kvůli tomu je při návrhu datových center potřeba zvážit několik faktorů:

\begin{description}
    \item[Výhody vodního chlazení] Systémy chlazení, kde je kapalina přiváděna pří\-mo k čipu (tzv. Direct-to-Chip), umožňují snížit spotřebu energie na chlazení až o 90~\%. Kapalinové chlazení je tedy velmi efektivní -- jak již bylo zmíněno, nejnovější datová centra používající tuto technologii dosahují PUE nižší než 1,1~\cite{sheng_power_2026, chen_electricity_2025}. Nejnovější systémy, jako je immersion cooling (servery jsou plně ponořené v kapalině), dosahují ještě vyšší efektivity (PUE < 1,03) a oproti vzduchovému chlazení mají o 95~\% nižší spotřebu energie.~\cite[s.~17--19]{cruzes_data_2025}
    \item[Nevýhody vodního chlazení] Vyšší účinnost je však vykoupena vysokými počátečními investicemi a vyšší technologickou náročností. Systémy jako D2C vyžadují specializované chladicí desky na procesorech, těsná potrubí uvnitř racků a jednotky pro distribuci kapaliny~\cite[s.~18]{cruzes_data_2025}. Navíc je obtížné přestavět vzduchem chlazená DC na vodní chlazení, neboť nemívají dostatečně vyztužené podlahy (kapalina tvoří další zátěž)~\cite[s.~11]{avelar_ai_2023}.
\end{description}

Ačkoli je tedy konstrukce vodního chlazení mnohem nákladnější, v dlouhodobém horizontu se tato investice může vrátit v nižších nákladech na elektřinu. Především se ale jedná o technologickou nutnost, která je způsobena zvyšující se hustotou výkonu na jeden rack. Nabízí se ale otázka, kde datová centra vodu na chlazení berou a zda to nepřináší další náklady -- finanční i environmentální.

\subsection{Odběr a spotřeba vody moderních DC}
\label{subsec:voda}

Přechod od vzduchového chlazení k pokročilým kapalinovým systémům dramaticky snížil spotřebu elektrické energie nutnou pro chlazení. Tento zisk v energetické efektivitě je však často na úkor zvýšených nároků na spotřebu vody. Při hodnocení vlivu DC na vodní zdroje je nutné rozlišovat mezi:

\begin{itemize}
    \item odběrem vody (water withdrawal), který představuje celkový objem sladké vody odčerpané z podzemních nebo povrchových zdrojů;
    \item a spotřebou vody (water consumption), která označuje nenávratně ztracenou část odebrané vody (nejčastěji se odpařuje do atmosféry);~\cite{li_making_2025}
    \item platí tedy vztah $\text{Water consumption} \le \text{Water withdrawal}.$
\end{itemize}

Datová centra přímo odebírají vodu pro své chladicí systémy (Scope 1). Teplo ze serverů se přenáší do vodního okruhu a v chladicí věži se část této vody cíleně odpaří, čímž se teplo uvolní do okolí. V okruhu by se nicméně mohly koncentrovat minerály, soli a bakterie, takže i nevypařená voda se musí pravidelně vypouštět a neustále doplňovat novou sladkou vodou. Z toho důvodu vyžadují některá vodou chlazená DC neustálý přítok vody, mnohdy pitné.~\cite{li_making_2025}

Mnohem větší zátěž na vodní zdroje však představuje tzv. nepřímý odběr (Scope 2). Jde o vodu odebíranou v elektrárnách, které pokrývají energetickou spotřebu datových center -- zejména tepelné či jaderné elektrárny potřebují k výrobě elektřiny masivní chlazení a vodní elektrárny zase ztrácejí vodu odparem z nádrží~\cite{li_making_2025, LBNL-2001637}. V USA je průměrný odběr vody pro výrobu elektřiny odhadován na 43,8 litru na každou spotřebovanou kilowatthodinu~\cite{li_making_2025}.

Pro představu, typické 100MW hyperscale datové centrum může spotřebovat kolem 2 milionů litrů vody denně (Scope 2), což odpovídá spotřebě zhruba 6 500 domácností. 60~\% z této spotřeby je nepřímá.~\cite{international_energy_agency_iea_energy_2025}


\section{Faktory pro výběr lokace AI datového centra}
\label{sec:lokalita}

%%%%%%%%%%%%%%%%%%%%%%%%%%%%%%%%%%%%%%%%%%%%%%%%%%%%%%%%%%%%%%%%%%%%%%%%%%%%%
% SEKCE 5: LOKALITA
% - Faktory ke zvážení při výstavbě datového centra
% - Formulováno jako PESTLE analýza pro lepší přehlednost, je to poměrně známý nástroj
%%%%%%%%%%%%%%%%%%%%%%%%%%%%%%%%%%%%%%%%%%%%%%%%%%%%%%%%%%%%%%%%%%%%%%%%%%%%%

Ken Baudry~\cite{geng_data_2014} uvádí, že výběr lokality funguje vylučovací metodou -- začíná se širokým geografickým záběrem a postupně se eliminují místa, která nesplňují kritické požadavky, až zbyde několik málo kandidátů. S nástupem AI datových center se navíc kritéria radikálně zpřísnila. Ta nejdůležitější byla autorem této práce rozdělena do skupin inspirovaných analýzou PESTLE:

\begin{description}
    \item[Politické faktory] Při mezinárodní expanzi podniky zohledňují zákony o och\-ra\-ně soukromí a geopolitická rizika spojená s tím, kde jsou data fyzicky uložena~\cite[s.~96]{geng_data_2014}. Důležitá je i stabilita podnikatelského prostředí, protože DC vyžadují obrovské investice. Z pohledu vlád se AI postupně stává záležitostí státní bezpečnosti, jednotlivé státy (nebo celky jako EU) se mohou snažit snížit svou závislost na zahraničních datových centrech, což může způsobit jejich širší distribuci po celém světě~\cite[s.~25]{pilz_compute_2023}. 

    \item[Ekonomické faktory] Náklady na výstavbu a provoz DC drasticky ovlivňují místní vlády, zejména na území USA. Provozovatelé proto hledají daňové pobídky a úlevy, například odpuštění daně z nemovitosti~\cite[s.~89--97]{geng_data_2014}. Moderní DC se navíc často stavějí dál od velkých metropolí -- v odlehlých oblastech lze rychleji sehnat stavební povolení a také v nich jdou stavět nové rozvodové sítě.~\cite[s.~25]{pilz_compute_2023}

    \item[Sociální faktory] Datová centra ke svému provozu vyžadují kvalifikovaný personál, kvůli čemuž je výhodou blízkost technických univerzit a výzkumných center~\cite{geng_data_2014}. Mnohem důležitějším aspektem jsou však iniciativy obyvatel -- stavba datových center vyvolává společenské konflikty, protože přímo soutěží o vodní zdroje s místními obyvateli a zemědělstvím~\cite[s.~4]{shah_four_2026}. Problémy představuje také nepřetržité hučení způsobené chladicími systémy a zhoršená dopravní infrastruktura.

    \item[Technologické faktory] Dostupnost spolehlivé a levné elektřiny je dnes hlavním faktorem. Protože výstavba nových přenosových soustav trvá 4 až 10 let, datová centra narážejí na přetížené sítě a dlouhé pořadníky pro připojení (fronty mohou trvat i roky).~\cite{ginzburg-ganz_technical_2025, international_energy_agency_iea_energy_2025} Z tohoto důvodu se nová DC často umisťují ke zdrojům energií, jak je popsáno v kap.~\ref{subsec:odkud_energie}. Pro inferenční centra je pak důležitá také konektivita (kap. \ref{subsec:ai_factory}), kvůli čemuž by měla být stavěna v blízkosti páteřní optické sítě. 

    \item[Legislativní faktory]
    Z pohledu legislativy je důležité, jakým způsobem je výstavba datových center regulována. V současnosti se již objevují specifické zákony, které se jejich výstavby a provozu přímo týkají. Například v Oregonu byl prosazen zákon, který datová centra zavazuje k alespoň 10letému provozu a zakazuje socializaci nákladů na infrastrukturu.~\cite[s.~19]{ginzburg-ganz_technical_2025}

    \item[Environmentální faktory] Kvůli chlazení se preferují lokality s chladnějším podnebím, které umožňují využívat venkovní vzduch po většinu roku~\cite{barroso_data_2026, geng_data_2014}. Pro odpařovací systémy je důležitá i dostupnost vody. Druhým environmentálním faktorem je stabilita a bezpečnost prostředí -- lokality se zkoumají z hlediska rizika zemětřesení, záplav, hurikánů, tornád či tsunami.~\cite{tatnall_reflections_2012,geng_data_2014}
\end{description}

\newpage

\section{Shrnutí}
\label{sec:shrnuti}

%%%%%%%%%%%%%%%%%%%%%%%%%%%%%%%%%%%%%%%%%%%%%%%%%%%%%%%%%%%%%%%%%%%%%%%%%%%%%%%%
% KAPITOLA 6: SHRNUTÍ
% - Stručně opakuje předchozí sekce
% - Je to to nejdůležitější, co je potřeba pro pochopení následujících kapitol
%%%%%%%%%%%%%%%%%%%%%%%%%%%%%%%%%%%%%%%%%%%%%%%%%%%%%%%%%%%%%%%%%%%%%%%%%%%%%%%%

Datová centra prošla za posledních 30 let radikální proměnou. Od sálových počítačů a nezávislých firemních serveroven v 90. letech se průmysl na počátku tisíciletí přesunul k tzv. WSC. Tyto obří komplexy propojily statisíce levných CPU serverů do jednoho fungujícího celku, aby obsloužily globální cloudové služby a vyhledávání. Nástup umělé inteligence však odstartoval novou éru -- WSC se mění na AI Factories, což jsou vysoce specializované superpočítače navržené nikoli pro běh mnoha nezávislých úloh, ale pro masivně paralelní trénink velkých jazykových modelů.~\cite{barroso_data_2026,barroso_brief_2021,geng_data_2014}

Tradiční servery postavené na univerzálních procesorech (CPU) přestaly pro AI stačit. Výpočetní zátěž převzaly specializované akcelerátory (jako GPU a TPU), které jsou optimalizované pro trénování AI. S tím však vzrostl tepelný výkon a energetická spotřeba samotných čipů (TDP) a nutnost umisťovat je fyzicky co nejblíže k sobě kvůli komunikační latenci. To vedlo k nárůstu výkonové hustoty racků (z tradičních 5--10~kW na 30--100+~kW).~\cite{barroso_data_2026, sheng_power_2026,chen_electricity_2025}

Kvůli tomu spotřeba energie datových center výrazně roste. AI datová centra vstupují do \uv{gigawattové éry} a zvyšují nároky na svou efektivitu (PUE). Zásadním problémem pro elektrickou síť ale není jen celkový objem, nýbrž pulzní chování AI zátěže -- synchronizované výpočty tisíců GPU generují obrovské výkyvy spotřeby v řádu desítek megawattů během milisekund, což ohrožuje frekvenční stabilitu sítě a zhoršuje kvalitu dodávek (harmonické zkreslení, problikávání). Jedno 200MW DC může ročně spotřebovat přibližně 1,4~TWh elektřiny, což je podobně jako spotřeba Karlovarského kraje.~\cite{sheng_power_2026,chen_electricity_2025, energeticky_regulacni_urad_spotreba_2024}

AI racky kvůli vysoké hustotě výkonu nelze spolehlivě chladit vzduchem. Nastupuje proto kapalinové chlazení, které je jak výkonnější, tak energeticky efektivnější. Na druhou stranu vyžaduje vysoké počáteční investice a spotřebuje obrovské množství vody (odpařuje se v chladicích věžích). Do spotřeby vody datovým centrem lze zahrnout také vodu, která byla spotřebována pro výrobu jím odebrané elektřiny. Běžné hyperscale centrum tak může denně spotřebovat až 2 miliony litrů vody, což je podobně jako 6 500 domácností.~\cite{jegham_how_2025,cruzes_data_2025,international_energy_agency_iea_energy_2025,shah_four_2026}

Při vybírání lokality pro výstavbu AI datového centra je hlavní prioritou kapacita a propustnost elektrické sítě (často se čeká roky na připojení). Hledají se místa s levnou a stabilní bezuhlíkovou energií (blízkost větrných, solárních či jaderných elektráren), dostatečným zdrojem vody pro chlazení, páteřní optickou sítí a příznivým legislativním i daňovým prostředím.~\cite{geng_data_2014, chen_electricity_2025, pilz_compute_2023, lin_exploding_2024}